% --- Archon physics formalization source begin ---
% archon:physics
% archon:covers PhyXMiniProblems/problem_phyx_mini_0245.lean
% archon:source-report reports/phyx_mini/problem_phyx_mini_0245.source.json

\chapter{Physics problem phyx\_mini\_0245}
\label{ch:PhyXMiniProblems_problem_phyx_mini_0245}

\paragraph{Problem source.}
\#\# Physical scenario

Helium-neon lasers emit the red laser light commonly used in classroom demonstrations and supermarket checkout scanners. A helium-neon laser operates at a wavelength of precisely \$632.9924 \textbackslash{}, \textbackslash{}mathrm\{nm\}\$ when the spacing between the mirrors is \$310.372 \textbackslash{}, \textbackslash{}mathrm\{mm\}\$.

\#\# Question

What is the next longest wavelength that could form a standing wave in this laser cavity?

\#\# Answer choices

- A: A: 632.9930nm

- B: B: 622.9930nm

- C: C: 642.5693nm

- D: D: 630.9130nm

\#\# Figure caption (auxiliary; use the image as primary evidence)

The image is a diagram depicting the basic structure of a laser. It includes the following elements and labels:

- **Laser cavity**: A horizontal rectangular area shown in light grey, where the laser action occurs.
- **Full reflector**: Located on the left side of the cavity, it is illustrated as a solid blue vertical line.
- **Partial reflector**: Found on the right side, depicted as a blue vertical line similar to the full reflector but allowing some light to pass through.
- **Standing light wave**: A purple wavy line inside the laser cavity, representing the oscillating light wave.
- **Laser beam**: A red horizontal line extending from the partial reflector on the right, indicating the emitted light.

The diagram visually explains how a laser cavity works, with a standing light wave bouncing between the full and partial reflectors, and a laser beam emerging from the partial reflector.

\#\# Dataset metadata

Resonance and Harmonics; Physical Model Grounding Reasoning, Spatial Relation Reasoning

\paragraph{Recorded answer/context.}
A: A: 632.9930nm

\paragraph{Figure/image path.}
/root/proposal\_for\_physic/science-mango/phyx\_mini\_run/phyx\_data/test\_image/245.png

\paragraph{Formalization target.}
create a compiling Lean file with sorry bodies at `PhyXMiniProblems/problem\_phyx\_mini\_0245.lean`. The Lean declarations must preserve the physical quantities, dimensions or dimensional roles, figure labels, governing-law hypotheses, and final relation expressed by this problem.
Use Mathlib/Physlib names found through LeanExplore where available. If a physics API is missing, introduce faithful local abstractions rather than scalar placeholder aliases.

\begin{theorem}[Physics formalization target]
\label{thm:physics:phyx_mini_0245:target}
\lean{PhyXMiniProblems.ProblemPhyXMini0245.problem_phyx_mini_0245}
\uses{def:physics:phyx-mini-0245:phyxminiproblems-problemphyxmini0245-lengthquantity, def:physics:phyx-mini-0245:phyxminiproblems-problemphyxmini0245-lengthinnanometers, def:physics:phyx-mini-0245:phyxminiproblems-problemphyxmini0245-roundstodisplayedvalue, def:physics:phyx-mini-0245:phyxminiproblems-problemphyxmini0245-heliumneonlasercavitysetup, def:physics:phyx-mini-0245:phyxminiproblems-problemphyxmini0245-matchesprimarylaserfigure, def:physics:phyx-mini-0245:phyxminiproblems-problemphyxmini0245-matchesproblemreadouts, def:physics:phyx-mini-0245:phyxminiproblems-problemphyxmini0245-hasphysicallaserparameters, def:physics:phyx-mini-0245:phyxminiproblems-problemphyxmini0245-satisfieslongitudinalcavityresonancelaw, def:physics:phyx-mini-0245:phyxminiproblems-problemphyxmini0245-isnextlongerstandingwavelength, def:physics:phyx-mini-0245:phyxminiproblems-problemphyxmini0245-answerchoice, def:physics:phyx-mini-0245:phyxminiproblems-problemphyxmini0245-matchesdisplayedanswer}
The assigned autoformalize agent should translate this physics problem into Lean declarations in the covered file, with theorem and lemma proof bodies written as `by sorry`.
\end{theorem}
\begin{proof}
This is an autoformalization task, not a proof task. Produce statements that can later be proved by the physics prover without weakening the physical model.
\end{proof}

% --- Archon declaration topology begin ---
\section{Declaration topology}
\begin{definition}[Length Quantity]
  \label{def:physics:phyx-mini-0245:phyxminiproblems-problemphyxmini0245-lengthquantity}
  \lean{PhyXMiniProblems.ProblemPhyXMini0245.LengthQuantity}
  A nonnegative physical length, independent of the unit used to read it
\end{definition}
\begin{proof}
  This is the declared model definition of Length Quantity.
\end{proof}
\begin{definition}[length Readout]
  \label{def:physics:phyx-mini-0245:phyxminiproblems-problemphyxmini0245-lengthreadout}
  \lean{PhyXMiniProblems.ProblemPhyXMini0245.lengthReadout}
  \uses{def:physics:phyx-mini-0245:phyxminiproblems-problemphyxmini0245-lengthquantity}
  Read a physical length as a real number in a selected length unit
\end{definition}
\begin{proof}
  This is the declared model definition of length Readout.
\end{proof}
\begin{definition}[length In Nanometers]
  \label{def:physics:phyx-mini-0245:phyxminiproblems-problemphyxmini0245-lengthinnanometers}
  \lean{PhyXMiniProblems.ProblemPhyXMini0245.lengthInNanometers}
  \uses{def:physics:phyx-mini-0245:phyxminiproblems-problemphyxmini0245-lengthquantity, def:physics:phyx-mini-0245:phyxminiproblems-problemphyxmini0245-lengthreadout}
  Nanometre readout, used for the laser wavelengths and answer choices
\end{definition}
\begin{proof}
  This is the declared model definition of length In Nanometers.
\end{proof}
\begin{definition}[length In Millimeters]
  \label{def:physics:phyx-mini-0245:phyxminiproblems-problemphyxmini0245-lengthinmillimeters}
  \lean{PhyXMiniProblems.ProblemPhyXMini0245.lengthInMillimeters}
  \uses{def:physics:phyx-mini-0245:phyxminiproblems-problemphyxmini0245-lengthquantity, def:physics:phyx-mini-0245:phyxminiproblems-problemphyxmini0245-lengthreadout}
  Millimetre readout, used for the stated mirror separation
\end{definition}
\begin{proof}
  This is the declared model definition of length In Millimeters.
\end{proof}
\begin{definition}[length In Meters]
  \label{def:physics:phyx-mini-0245:phyxminiproblems-problemphyxmini0245-lengthinmeters}
  \lean{PhyXMiniProblems.ProblemPhyXMini0245.lengthInMeters}
  \uses{def:physics:phyx-mini-0245:phyxminiproblems-problemphyxmini0245-lengthquantity, def:physics:phyx-mini-0245:phyxminiproblems-problemphyxmini0245-lengthreadout}
  Metre readout, used to state positivity independently of printed scales
\end{definition}
\begin{proof}
  This is the declared model definition of length In Meters.
\end{proof}
\begin{definition}[Rounds To Displayed Value]
  \label{def:physics:phyx-mini-0245:phyxminiproblems-problemphyxmini0245-roundstodisplayedvalue}
  \lean{PhyXMiniProblems.ProblemPhyXMini0245.RoundsToDisplayedValue}
  RoundsToDisplayedValue actual displayed resolution means that actual is in the half-open rounding bin centered on displayed. The half-open choice fixes the usual tie convention without affecting this problem's values
\end{definition}
\begin{proof}
  This is the declared model definition of Rounds To Displayed Value.
\end{proof}
\begin{definition}[Laser Gain Medium]
  \label{def:physics:phyx-mini-0245:phyxminiproblems-problemphyxmini0245-lasergainmedium}
  \lean{PhyXMiniProblems.ProblemPhyXMini0245.LaserGainMedium}
  The gain medium named in the problem statement
\end{definition}
\begin{proof}
  This is the declared model definition of Laser Gain Medium.
\end{proof}
\begin{definition}[Laser Light Color]
  \label{def:physics:phyx-mini-0245:phyxminiproblems-problemphyxmini0245-laserlightcolor}
  \lean{PhyXMiniProblems.ProblemPhyXMini0245.LaserLightColor}
  The visible color of the emitted laser light
\end{definition}
\begin{proof}
  This is the declared model definition of Laser Light Color.
\end{proof}
\begin{definition}[Cavity End]
  \label{def:physics:phyx-mini-0245:phyxminiproblems-problemphyxmini0245-cavityend}
  \lean{PhyXMiniProblems.ProblemPhyXMini0245.CavityEnd}
  The two ends of the horizontal laser cavity in the primary figure
\end{definition}
\begin{proof}
  This is the declared model definition of Cavity End.
\end{proof}
\begin{definition}[Reflector Kind]
  \label{def:physics:phyx-mini-0245:phyxminiproblems-problemphyxmini0245-reflectorkind}
  \lean{PhyXMiniProblems.ProblemPhyXMini0245.ReflectorKind}
  Optical behavior of a cavity reflector
\end{definition}
\begin{proof}
  This is the declared model definition of Reflector Kind.
\end{proof}
\begin{definition}[Figure Element]
  \label{def:physics:phyx-mini-0245:phyxminiproblems-problemphyxmini0245-figureelement}
  \lean{PhyXMiniProblems.ProblemPhyXMini0245.FigureElement}
  The five labeled physical elements visible in the primary figure
\end{definition}
\begin{proof}
  This is the declared model definition of Figure Element.
\end{proof}
\begin{definition}[Figure Placement]
  \label{def:physics:phyx-mini-0245:phyxminiproblems-problemphyxmini0245-figureplacement}
  \lean{PhyXMiniProblems.ProblemPhyXMini0245.FigurePlacement}
  Qualitative horizontal placement of a labeled element in the figure
\end{definition}
\begin{proof}
  This is the declared model definition of Figure Placement.
\end{proof}
\begin{definition}[Helium Neon Laser Cavity Setup]
  \label{def:physics:phyx-mini-0245:phyxminiproblems-problemphyxmini0245-heliumneonlasercavitysetup}
  \lean{PhyXMiniProblems.ProblemPhyXMini0245.HeliumNeonLaserCavitySetup}
  \uses{def:physics:phyx-mini-0245:phyxminiproblems-problemphyxmini0245-lengthquantity, def:physics:phyx-mini-0245:phyxminiproblems-problemphyxmini0245-lasergainmedium, def:physics:phyx-mini-0245:phyxminiproblems-problemphyxmini0245-laserlightcolor, def:physics:phyx-mini-0245:phyxminiproblems-problemphyxmini0245-cavityend, def:physics:phyx-mini-0245:phyxminiproblems-problemphyxmini0245-reflectorkind, def:physics:phyx-mini-0245:phyxminiproblems-problemphyxmini0245-figureelement, def:physics:phyx-mini-0245:phyxminiproblems-problemphyxmini0245-figureplacement}
  Independent physical quantities and qualitative attributes of the laser. formsLongitudinalStandingWave is an observational predicate; its governing relation to mode number and mirror separation is stated separately below. No next-longer wavelength or answer choice is stored in this setup
\end{definition}
\begin{proof}
  This is the declared model definition of Helium Neon Laser Cavity Setup.
\end{proof}
\begin{definition}[Matches Primary Laser Figure]
  \label{def:physics:phyx-mini-0245:phyxminiproblems-problemphyxmini0245-matchesprimarylaserfigure}
  \lean{PhyXMiniProblems.ProblemPhyXMini0245.MatchesPrimaryLaserFigure}
  \uses{def:physics:phyx-mini-0245:phyxminiproblems-problemphyxmini0245-figureelement, def:physics:phyx-mini-0245:phyxminiproblems-problemphyxmini0245-heliumneonlasercavitysetup}
  Information read from the primary image. The purple standing wave is inside the cavity, between a full left reflector and a partial right reflector; the red output beam extends to the right of the partial reflector
\end{definition}
\begin{proof}
  This is the declared model definition of Matches Primary Laser Figure.
\end{proof}
\begin{definition}[Matches Problem Readouts]
  \label{def:physics:phyx-mini-0245:phyxminiproblems-problemphyxmini0245-matchesproblemreadouts}
  \lean{PhyXMiniProblems.ProblemPhyXMini0245.MatchesProblemReadouts}
  \uses{def:physics:phyx-mini-0245:phyxminiproblems-problemphyxmini0245-lengthinnanometers, def:physics:phyx-mini-0245:phyxminiproblems-problemphyxmini0245-lengthinmillimeters, def:physics:phyx-mini-0245:phyxminiproblems-problemphyxmini0245-roundstodisplayedvalue, def:physics:phyx-mini-0245:phyxminiproblems-problemphyxmini0245-heliumneonlasercavitysetup}
  Problem-statement data. The wavelength marked "precisely" is an exact nanometre readout. The mirror spacing is placed in the rounding bin denoted by the printed 310.372 mm, so the data remain compatible with an integral longitudinal mode. Neither the next wavelength nor answer A occurs here
\end{definition}
\begin{proof}
  This is the declared model definition of Matches Problem Readouts.
\end{proof}
\begin{definition}[Has Physical Laser Parameters]
  \label{def:physics:phyx-mini-0245:phyxminiproblems-problemphyxmini0245-hasphysicallaserparameters}
  \lean{PhyXMiniProblems.ProblemPhyXMini0245.HasPhysicalLaserParameters}
  \uses{def:physics:phyx-mini-0245:phyxminiproblems-problemphyxmini0245-lengthinmeters, def:physics:phyx-mini-0245:phyxminiproblems-problemphyxmini0245-heliumneonlasercavitysetup}
  Positivity conditions for the nondegenerate optical cavity and wave
\end{definition}
\begin{proof}
  This is the declared model definition of Has Physical Laser Parameters.
\end{proof}
\begin{definition}[Satisfies Longitudinal Cavity Resonance Law]
  \label{def:physics:phyx-mini-0245:phyxminiproblems-problemphyxmini0245-satisfieslongitudinalcavityresonancelaw}
  \lean{PhyXMiniProblems.ProblemPhyXMini0245.SatisfiesLongitudinalCavityResonanceLaw}
  \uses{def:physics:phyx-mini-0245:phyxminiproblems-problemphyxmini0245-lengthquantity, def:physics:phyx-mini-0245:phyxminiproblems-problemphyxmini0245-lengthreadout, def:physics:phyx-mini-0245:phyxminiproblems-problemphyxmini0245-heliumneonlasercavitysetup}
  The longitudinal cavity-resonance law. A wavelength forms a standing wave exactly when a positive integral number of half-wavelengths fits between the reflectors: 2 L = n lambda. It is stated in every length unit and for every candidate wavelength, so it is a general governing law rather than the requested adjacent-mode formula
\end{definition}
\begin{proof}
  This is the declared model definition of Satisfies Longitudinal Cavity Resonance Law.
\end{proof}
\begin{definition}[Is Next Longer Standing Wavelength]
  \label{def:physics:phyx-mini-0245:phyxminiproblems-problemphyxmini0245-isnextlongerstandingwavelength}
  \lean{PhyXMiniProblems.ProblemPhyXMini0245.IsNextLongerStandingWavelength}
  \uses{def:physics:phyx-mini-0245:phyxminiproblems-problemphyxmini0245-lengthquantity, def:physics:phyx-mini-0245:phyxminiproblems-problemphyxmini0245-lengthinnanometers, def:physics:phyx-mini-0245:phyxminiproblems-problemphyxmini0245-heliumneonlasercavitysetup}
  A candidate is the next longer standing wavelength when it is resonant, strictly longer than the operating wavelength, and no other resonant wavelength lies strictly between them. This definition characterizes the question being asked; it contains no numerical candidate or answer value
\end{definition}
\begin{proof}
  This is the declared model definition of Is Next Longer Standing Wavelength.
\end{proof}
\begin{definition}[Answer Choice]
  \label{def:physics:phyx-mini-0245:phyxminiproblems-problemphyxmini0245-answerchoice}
  \lean{PhyXMiniProblems.ProblemPhyXMini0245.AnswerChoice}
  Labels attached to the four wavelength choices in the source problem
\end{definition}
\begin{proof}
  This is the declared model definition of Answer Choice.
\end{proof}
\begin{definition}[displayed Answer Wavelength In Nanometers]
  \label{def:physics:phyx-mini-0245:phyxminiproblems-problemphyxmini0245-displayedanswerwavelengthinnanometers}
  \lean{PhyXMiniProblems.ProblemPhyXMini0245.displayedAnswerWavelengthInNanometers}
  \uses{def:physics:phyx-mini-0245:phyxminiproblems-problemphyxmini0245-answerchoice}
  Nanometre values printed beside the four answer labels
\end{definition}
\begin{proof}
  This is the declared model definition of displayed Answer Wavelength In Nanometers.
\end{proof}
\begin{definition}[recorded Dataset Answer]
  \label{def:physics:phyx-mini-0245:phyxminiproblems-problemphyxmini0245-recordeddatasetanswer}
  \lean{PhyXMiniProblems.ProblemPhyXMini0245.recordedDatasetAnswer}
  \uses{def:physics:phyx-mini-0245:phyxminiproblems-problemphyxmini0245-answerchoice}
  Dataset answer metadata, recorded as data rather than used as a premise
\end{definition}
\begin{proof}
  This is the declared model definition of recorded Dataset Answer.
\end{proof}
\begin{definition}[Matches Displayed Answer]
  \label{def:physics:phyx-mini-0245:phyxminiproblems-problemphyxmini0245-matchesdisplayedanswer}
  \lean{PhyXMiniProblems.ProblemPhyXMini0245.MatchesDisplayedAnswer}
  \uses{def:physics:phyx-mini-0245:phyxminiproblems-problemphyxmini0245-lengthquantity, def:physics:phyx-mini-0245:phyxminiproblems-problemphyxmini0245-lengthinnanometers, def:physics:phyx-mini-0245:phyxminiproblems-problemphyxmini0245-roundstodisplayedvalue, def:physics:phyx-mini-0245:phyxminiproblems-problemphyxmini0245-answerchoice, def:physics:phyx-mini-0245:phyxminiproblems-problemphyxmini0245-displayedanswerwavelengthinnanometers}
  A physical wavelength rounds to the value printed for an answer choice
\end{definition}
\begin{proof}
  This is the declared model definition of Matches Displayed Answer.
\end{proof}
\begin{lemma}[operating Mode Number is near 980650]
  \label{lem:physics:phyx-mini-0245:phyxminiproblems-problemphyxmini0245-operatingmodenumber-is-near-980650}
  \lean{PhyXMiniProblems.ProblemPhyXMini0245.operatingModeNumber_is_near_980650}
  \uses{def:physics:phyx-mini-0245:phyxminiproblems-problemphyxmini0245-lengthreadout, def:physics:phyx-mini-0245:phyxminiproblems-problemphyxmini0245-heliumneonlasercavitysetup, def:physics:phyx-mini-0245:phyxminiproblems-problemphyxmini0245-matchesproblemreadouts, def:physics:phyx-mini-0245:phyxminiproblems-problemphyxmini0245-satisfieslongitudinalcavityresonancelaw}
  The exact operating wavelength and rounded mirror spacing restrict the current longitudinal mode to one of three adjacent large integers. This intermediate statement records the arithmetic route without assuming the next wavelength
\end{lemma}
\begin{proof}
  The conclusion follows from the governing relations and figure readouts named in the statement of operating Mode Number is near 980650.
\end{proof}
% --- Archon declaration topology end ---
% --- Archon physics formalization source end ---
